\documentclass[12pt]{article}
\usepackage[utf8]{inputenc}
\usepackage[T1]{fontenc}
\usepackage[margin=1in]{geometry}
\usepackage{setspace}
\usepackage{xcolor}
\usepackage[colorlinks=true, urlcolor=blue!60!black, linkcolor=blue!60!black]{hyperref}
\setlength{\parindent}{2em}
\setlength{\parskip}{1.25em}
\renewcommand{\baselinestretch}{1.0}
\begin{document}
\begin{center}
{\LARGE \textit{Youngstown Sheet \& Tube Co. v. Sawyer}, 343 U.S. 579 (1952)\par}
\vspace{0.8em}
{\large Case Brief\par}
\end{center}
\vspace{1em}

\section{Memory Jogger}
In the landmark ``Steel Seizure Case,'' a 6--3 Supreme Court held that President Truman lacked both statutory and constitutional authority to seize the Nation's strike-bound steel mills during the Korean War, establishing Justice Jackson's celebrated three-zone framework for judging executive power (\href{https://www.courtlistener.com/opinion/105018/youngstown-sheet-tube-co-v-sawyer/}{CourtListener}; \href{https://supreme.justia.com/cases/federal/us/343/579/}{Justia}).

\section{Detailed Case Facts}
In June 1950, following North Korea's invasion of South Korea and a United Nations Security Council resolution, President Harry S. Truman committed U.S. armed forces to the Korean War without a formal congressional declaration of war (\href{https://supreme.justia.com/cases/federal/us/343/579/}{Justia}). In December 1950, Truman proclaimed a national emergency, and in late 1951 a bitter collective-bargaining dispute erupted between the Nation's major steel companies---including Youngstown Sheet \& Tube Co., Bethlehem Steel, and Republic Steel---and their employees, represented by the United Steelworkers of America, CIO (\href{https://supreme.justia.com/cases/federal/us/343/579/}{Justia}). When negotiations stalled and the union announced a strike for December 31, 1951, Truman referred the dispute on December 22, 1951, to the Federal Wage Stabilization Board (WSB) established under Executive Order 10233 (\href{https://supreme.justia.com/cases/federal/us/343/579/}{Justia}).

The WSB's wage-increase recommendations failed to produce a settlement, and on April 4, 1952, the union gave notice of a nationwide strike set to begin at 12:01 a.m. on April 9, 1952 (\href{https://supreme.justia.com/cases/federal/us/343/579/}{Justia}). Believing that an uninterrupted supply of steel was indispensable to weapons production and atomic-energy defense programs, Truman declined to invoke the 80-day cooling-off injunction procedure of the Labor Management Relations Act of 1947 (the Taft-Hartley Act), which he had vetoed and considered unfair to labor, and declined to use the slower condemnation procedures of the Selective Service Act of 1948 or the Defense Production Act of 1950 (\href{https://supreme.justia.com/cases/federal/us/343/579/}{Justia}; \href{https://www.law.cornell.edu/wex/youngstown_sheet_tube_co._v._sawyer_\%281952\%29}{Cornell LII Wex}).

Instead, on the evening of April 8, 1952---hours before the strike deadline---Truman issued Executive Order 10340, directing Secretary of Commerce Charles Sawyer to take possession of and operate most of the Nation's steel mills (\href{https://supreme.justia.com/cases/federal/us/343/579/}{Justia}). The Order cited no specific statute, invoking instead ``the authority vested in me by the Constitution and laws of the United States, and as President of the United States and Commander in Chief of the armed forces'' (\href{https://supreme.justia.com/cases/federal/us/343/579/}{Justia}). Sawyer immediately issued possessory orders designating the steel companies' presidents as operating managers for the United States. Truman sent messages reporting the seizure to Congress on April 9 and April 21, 1952, inviting legislative action, but Congress took no action (\href{https://supreme.justia.com/cases/federal/us/343/579/}{Justia}). Obeying under protest, the steel companies sued Secretary Sawyer in federal district court seeking declaratory and injunctive relief (\href{https://supreme.justia.com/cases/federal/us/343/579/}{Justia}).

\section{Procedural History}
After an initial request for a temporary restraining order was denied by Judge Alexander Holtzoff, the steel companies' motion for a preliminary injunction came before Judge David A. Pine of the U.S. District Court for the District of Columbia (\href{https://supreme.justia.com/cases/federal/us/343/579/}{Justia}). On April 29--30, 1952, Judge Pine rejected the government's sweeping claim of unreviewable ``inherent'' emergency power and issued a preliminary injunction restraining Secretary Sawyer from continuing the seizure (\textit{Youngstown Sheet \& Tube Co. v. Sawyer}, 103 F. Supp. 569 (D.D.C. 1952)) (\href{https://supreme.justia.com/cases/federal/us/343/579/}{Justia}).

The U.S. Court of Appeals for the District of Columbia Circuit stayed Judge Pine's injunction on the same day, April 30, 1952 (197 F.2d 582 (D.C. Cir. 1952)) (\href{https://supreme.justia.com/cases/federal/us/343/579/}{Justia}). Bypassing full appellate review due to the national urgency, the U.S. Supreme Court granted certiorari on May 3, 1952 (343 U.S. 937), heard oral argument on May 12--13, 1952, and issued its decision just three weeks later on June 2, 1952, affirming the district court's injunction (\href{https://supreme.justia.com/cases/federal/us/343/579/}{Justia}; \href{https://www.law.cornell.edu/supremecourt/text/343/579}{Cornell LII}).

\section{Judicial Votes}
\textbf{Vote:} 6--3, affirming the District Court's preliminary injunction (\href{https://supreme.justia.com/cases/federal/us/343/579/}{Justia}). Every Justice in the six-member majority authored a separate opinion.
\begin{itemize}
\item \textbf{Majority Opinion (5):} Justice Black delivered the opinion of the Court, joined by Justices Frankfurter, Douglas, Jackson, and Burton.
\item \textbf{Concurrences (4 joining Black's opinion):} Justices Frankfurter, Douglas, Jackson, and Burton each filed separate concurring opinions.
\item \textbf{Concurrence in the Judgment (1):} Justice Clark concurred in the judgment of the Court only, without joining Justice Black's opinion.
\item \textbf{Dissent (3):} Chief Justice Vinson filed a dissenting opinion, joined by Justices Reed and Minton.
\end{itemize}

\section{Holding}
\textbf{Judgment:} Affirmed. Executive Order 10340 was unconstitutional, and the District Court's preliminary injunction barring Secretary Sawyer from seizing and operating the steel mills was upheld (\href{https://supreme.justia.com/cases/federal/us/343/579/}{Justia}).

\textbf{Rule Announced:} The President's power to issue an executive order seizing private domestic property must stem either from an act of Congress or from the Constitution itself. Neither the Commander-in-Chief Clause, the Executive Vesting Clause, nor the Take Care Clause of Article II authorizes the President to seize private industrial facilities to avert a domestic labor strike---a lawmaking function vested exclusively in Congress under Article I, especially where Congress has considered and declined to grant seizure power (\href{https://supreme.justia.com/cases/federal/us/343/579/}{Justia}; \href{https://www.law.cornell.edu/wex/youngstown_sheet_tube_co._v._sawyer_\%281952\%29}{Cornell LII Wex}).

\textbf{Holding vs. Dicta:}
\begin{itemize}
\item \textit{Binding Holding:} (1) The constitutional challenge was ripe for final resolution at the preliminary-injunction stage because seizure of going businesses threatened unquantifiable irreparable harm and recovery in the Court of Claims was doubtful; and (2) President Truman lacked statutory or Article II constitutional authority to seize the Nation's steel mills to settle a domestic labor dispute (\href{https://supreme.justia.com/cases/federal/us/343/579/}{Justia}).
\item \textit{Dicta and Doctrinal Evolution:} Justice Black's rigid formalist suggestion that the President never possesses lawmaking or emergency authority absent an explicit statutory or constitutional grant was tempered by the four concurrences joining his opinion. Crucially, Justice Jackson's solo concurring tripartite framework---technically non-binding dicta in 1952---subsequently became the Court's governing doctrinal test for separation-of-powers disputes in cases such as \textit{Dames \& Moore v. Regan}, 453 U.S. 654 (1981), \textit{Medell\'{i}n v. Texas}, 552 U.S. 491 (2008), and \textit{Zivotofsky v. Kerry}, 576 U.S. 1 (2015) (\href{https://www.law.cornell.edu/wex/youngstown_sheet_tube_co._v._sawyer_\%281952\%29}{Cornell LII Wex}; \href{https://www.uclalawreview.org/wp-content/uploads/2019/09/44_54UCLALRev1703August2007.pdf}{UCLA Law Review}).
\end{itemize}

\section{Analysis of Opinions}

\subsection{Majority Opinion (Justice Black)}
Justice Black wrote a brisk, strictly formalist and textualist opinion for five Justices (\href{https://constitutioncenter.org/blog/the-youngstown-case-three-approaches-to-interpreting-presidential-power}{National Constitution Center}). After holding that the case was ripe for constitutional adjudication despite its preliminary posture, Black announced a categorical rule: ``The President's power, if any, to issue the order must stem either from an act of Congress or from the Constitution itself'' (\href{https://supreme.justia.com/cases/federal/us/343/579/}{Justia}). No statute authorized the seizure; indeed, when enacting the Taft-Hartley Act in 1947, Congress had expressly rejected an amendment granting emergency seizure power to resolve labor disputes. Turning to Article II, Black rejected the government's claim of ``aggregate'' executive authority: the Commander-in-Chief Clause does not extend ``theater of war'' military powers to seizing domestic factories to resolve labor disputes, and the Take Care Clause ``refutes the idea that [the President] is to be a lawmaker,'' confining him to recommending and vetoing legislation. Because Executive Order 10340 proclaimed a presidential policy and directed its execution like a statute, it usurped the legislative power vested exclusively in Congress under Article I, Section 1 and the Necessary and Proper Clause (\href{https://supreme.justia.com/cases/federal/us/343/579/}{Justia}).

\subsection{Concurring Opinion (Justice Frankfurter)}
While joining Black's opinion, Justice Frankfurter argued that separation-of-powers boundaries are ``more complicated and flexible'' than Black's bright-line formalism suggested (\href{https://supreme.justia.com/cases/federal/us/343/579/}{Justia}). Frankfurter conducted an exhaustive review of 16 federal seizure statutes since 1916 (appended to his opinion) and the legislative history of the 1947 Taft-Hartley Act and 1950 Defense Production Act, demonstrating that Congress had consciously withheld general seizure power and required the President to report deadlocks for ad hoc legislation. Crucially, Frankfurter articulated his influential ``historical gloss'' doctrine: ``a systematic, unbroken, executive practice, long pursued to the knowledge of the Congress and never before questioned,'' may serve as a constitutional gloss on ``executive Power'' under Article II (\href{https://supreme.justia.com/cases/federal/us/343/579/}{Justia}; \href{https://constitutioncenter.org/blog/the-youngstown-case-three-approaches-to-interpreting-presidential-power}{National Constitution Center}). Surveying past industrial seizures, however, he found only three isolated pre-Pearl Harbor Roosevelt seizures without statutory or declared-war authority---far too sparse to establish congressional acquiescence comparable to \textit{United States v. Midwest Oil Co.}, 236 U.S. 459 (1915).

\subsection{Concurring Opinion (Justice Douglas)}
Justice Douglas joined Black's opinion and emphasized both structural formalism and the Fifth Amendment Takings Clause (\href{https://supreme.justia.com/cases/federal/us/343/579/}{Justia}). He observed that an emergency ``did not create power; it merely marked an occasion when power should be exercised,'' and quoted Justice Brandeis's famous warning in \textit{Myers v. United States}, 272 U.S. 52, 293 (1926), that separation of powers was adopted ``not to promote efficiency, but to preclude the exercise of arbitrary power.'' Douglas added a distinctive constitutional argument: seizing an industrial plant is a Fifth Amendment taking requiring ``just compensation,'' and because only Congress possesses the Article I, Section 8 power to raise revenues and appropriate funds to pay compensation, only Congress can authorize a lawful condemnation (\href{https://supreme.justia.com/cases/federal/us/343/579/}{Justia}).

\subsection{Concurring Opinion (Justice Jackson)}
Drawing on his prior experience as Attorney General to President Franklin D. Roosevelt, Justice Jackson offered a pragmatic, structural, and functionalist framework that has become the most influential opinion in the case (\href{https://supreme.justia.com/cases/federal/us/343/579/}{Justia}; \href{https://constitutioncenter.org/blog/the-youngstown-case-three-approaches-to-interpreting-presidential-power}{National Constitution Center}). Rejecting isolated textualism, Jackson observed that ``[p]residential powers are not fixed but fluctuate depending upon their disjunction or conjunction with those of Congress,'' and divided executive action into three famous categories:
\begin{enumerate}
\item \textit{Zone 1 (Maximum Authority):} When the President acts pursuant to an express or implied congressional authorization, his authority is ``at its maximum,'' combining his own Article II powers with all that Congress can delegate.
\item \textit{Zone 2 (Zone of Twilight):} When the President acts in the absence of a congressional grant or denial, he relies solely on his independent powers, operating in a ``zone of twilight'' where concurrent authority is uncertain and outcomes depend on ``the imperatives of events and contemporary imponderables.''
\item \textit{Zone 3 (Lowest Ebb):} When the President takes measures ``incompatible with the expressed or implied will of Congress,'' his power is ``at its lowest ebb,'' sustainable only if his exclusive Article II authority disables Congress from acting on the subject (\href{https://supreme.justia.com/cases/federal/us/343/579/}{Justia}).
\end{enumerate}
Because Congress had covered property seizures through three statutory schemes (the Selective Service Act, Defense Production Act, and Taft-Hartley Act) inconsistent with Truman's unilateral order, Jackson placed the seizure in Zone 3. Rejecting the Solicitor General's sweeping readings of the Vesting, Commander-in-Chief, and Take Care Clauses---and warning via the Weimar Republic's collapse against unrestrained executive emergency powers---Jackson concluded that the Commander in Chief of the Army and Navy is not ``Commander in Chief of the country, its industries and its inhabitants'' (\href{https://supreme.justia.com/cases/federal/us/343/579/}{Justia}).

\subsection{Concurring Opinion (Justice Burton)}
Justice Burton joined Black's opinion while emphasizing the statutory conflict created by Congress's deliberate choices in the Taft-Hartley Act (\href{https://supreme.justia.com/cases/federal/us/343/579/}{Justia}). Burton noted that Congress had prescribed specific, non-seizure procedures for national-emergency labor disputes and expressly reserved the option of emergency seizure for future legislative action. Without deciding what residual authority a President might exercise in an imminent military invasion or catastrophic domestic emergency where Congress had not spoken, Burton concluded that Truman could not bypass Congress's clear statutory framework (\href{https://supreme.justia.com/cases/federal/us/343/579/}{Justia}).

\subsection{Concurring in the Judgment (Justice Clark)}
Justice Clark---who had previously served as Truman's Attorney General---declined to join Black's opinion and concurred only in the judgment (\href{https://supreme.justia.com/cases/federal/us/343/579/}{Justia}; \href{https://www.roberthjackson.org/wp-content/uploads/migrated-files/article-calendar-files/korman_primer.pdf}{Robert H. Jackson Center}). Unlike Black and Douglas, Clark believed the Constitution grants the President broad independent authority to act in grave emergencies when Congress has not spoken, citing \textit{In re Neagle}, 135 U.S. 1 (1890), and \textit{Little v. Barreme}, 6 U.S. (2 Cranch) 170 (1804). However, where Congress has laid down specific statutory procedures to address an emergency---as it did in the Taft-Hartley Act, Selective Service Act, and Defense Production Act---the President must follow those statutory pathways rather than inventing a unilateral seizure remedy (\href{https://supreme.justia.com/cases/federal/us/343/579/}{Justia}).

\subsection{Dissenting Opinion (Chief Justice Vinson)}
Chief Justice Vinson, joined by Justices Reed and Minton, delivered a vigorous defense of strong executive crisis leadership (\href{https://supreme.justia.com/cases/federal/us/343/579/}{Justia}; \href{https://static.c-span.org/landmarkCases/pdf/Youngstown_Vinson_Dissent.pdf}{C-SPAN Vinson Dissent}). Vinson stressed the gravity of the Korean conflict and cited extensive historical precedents of unilateral presidential emergency action---from Washington, Jefferson, and Lincoln through Wilson and Franklin D. Roosevelt---to argue that the President's duty to ``take Care that the Laws be faithfully executed'' included preserving vital defense-procurement and military-aid statutes from collapse during a nationwide steel shutdown. Vinson insisted that Truman had not usurped legislative power, because he promptly reported the temporary seizure to Congress and pledged to abide by any congressional directive. Mocking the majority's view as a ``messenger-boy concept of the Office,'' Vinson argued that the Constitution did not render the Chief Executive helpless to preserve the status quo until Congress could act (\href{https://supreme.justia.com/cases/federal/us/343/579/}{Justia}; \href{https://www.roberthjackson.org/wp-content/uploads/migrated-files/article-calendar-files/korman_primer.pdf}{Robert H. Jackson Center}).

\section{Examples: Future Applications}
\begin{enumerate}
\item \textbf{Same-Side (Unconstitutional):} \textit{Unilateral Presidential Nationalization of Semiconductor Plants.} Amid an undeclared overseas crisis, the President issues an executive order seizing domestic semiconductor factories to avert a strike, bypassing an existing federal labor statute in which Congress expressly rejected plant-seizure remedies. \textit{Why:} Directly parallel to \textit{Youngstown}, this domestic property seizure falls into Jackson's Zone 3 (and fails Black's formalist test) because it contradicts Congress's statutory framework and lies outside exclusive Article II military command (\href{https://supreme.justia.com/cases/federal/us/343/579/}{Justia}).

\item \textbf{Same-Side (Unconstitutional):} \textit{Domestic Surveillance Violating a Comprehensive Federal Statute.} Claiming inherent Commander-in-Chief authority to gather intelligence against foreign threats, the President authorizes warrantless domestic wiretapping of U.S. citizens in direct violation of the Foreign Intelligence Surveillance Act (FISA). \textit{Why:} Because Congress has validly regulated domestic surveillance under Article I, the President's power is at its ``lowest ebb'' in Zone 3 and cannot override a binding statutory prohibition (\href{https://www.uclalawreview.org/wp-content/uploads/2019/09/44_54UCLALRev1703August2007.pdf}{UCLA Law Review}).

\item \textbf{Opposite-Side (Constitutional):} \textit{Emergency Asset Freezes Pursuant to Congressional Authorization.} Relying on the International Emergency Economic Powers Act (IEEPA), the President freezes foreign sovereign assets and nullifies domestic legal attachments as part of an international diplomatic settlement, with a long history of congressional acquiescence. \textit{Why:} As in \textit{Dames \& Moore v. Regan}, 453 U.S. 654 (1981), express or implied statutory authorization combined with historical practice places the President in Jackson's Zone 1 (or the upper reach of Zone 2), where executive authority is at its maximum (\href{https://www.uclalawreview.org/wp-content/uploads/2019/09/44_54UCLALRev1703August2007.pdf}{UCLA Law Review}).

\item \textbf{Opposite-Side (Constitutional):} \textit{Exclusive Presidential Recognition of Foreign Sovereigns.} Congress passes a statute requiring the State Department to list a disputed city as part of a specific foreign nation on U.S. passports, contrary to the President's formal foreign-policy recognition decision. \textit{Why:} As the Court held in \textit{Zivotofsky v. Kerry}, 576 U.S. 1 (2015), even in Jackson's Zone 3, the President prevails when exercising an \textit{exclusive} Article II power---such as the Reception Clause recognition power---that disables Congress from acting on the subject (\href{https://www.law.cornell.edu/wex/youngstown_sheet_tube_co._v._sawyer_\%281952\%29}{Cornell LII Wex}).

\item \textbf{Fence-Sitter (Close Call):} \textit{Temporary Cybersecurity Directive to Private Power Utilities Amid Congressional Silence.} Facing an imminent foreign state-sponsored cyberattack on the domestic electrical grid, the President issues a 30-day emergency executive order mandating specific defensive software configurations at privately owned utilities; Congress has neither authorized nor forbidden such mandatory directives. \textit{Why:} Under Black's strict formalism, the order looks like unauthorized domestic lawmaking; under Jackson's Zone 2 (``zone of twilight'') and Clark's concurrence, however, congressional silence combined with an urgent external security threat creates a classic constitutional borderline dependent on ``the imperatives of events and contemporary imponderables'' (\href{https://supreme.justia.com/cases/federal/us/343/579/}{Justia}).
\end{enumerate}

\section{Critique}
\textbf{Formalist and Originalist Perspectives:} Formalist and textualist scholars often praise Justice Black's majority opinion as a clear, principled defense of Article I's vesting of all federal lawmaking power in Congress, preventing the executive from eroding constitutional boundaries through ad hoc necessity claims (\href{https://constitutioncenter.org/blog/the-youngstown-case-three-approaches-to-interpreting-presidential-power}{National Constitution Center}; \href{https://openyls.law.yale.edu/server/api/core/bitstreams/6d79e130-b40f-4083-bbb4-5f93642bb432/content}{Yale Law Scholarship}). At the same time, unitary-executive originalists and scholars such as Jack Goldsmith and John Manning have argued that Black's opinion overstated the rigidity of Article II, noting that modern administrative practice, prosecutorial discretion, and foreign-affairs initiatives partially vindicate Chief Justice Vinson's concept of a presidential ``completion power'' to carry existing legislative schemes into effect (\href{https://openyls.law.yale.edu/bitstreams/d907a40e-541f-45f4-b625-3aa1d23c542b/download}{Goldsmith \& Manning, Yale L.J.}).

\textbf{Functionalist, Progressive, and Doctrinal Perspectives:} Functionalist and structural scholars widely celebrate Justice Jackson's concurrence for capturing the pragmatic interdependence of the political branches rather than treating Articles I and II as hermetically sealed compartments (\href{https://constitutioncenter.org/blog/the-youngstown-case-three-approaches-to-interpreting-presidential-power}{National Constitution Center}). Nevertheless, legal scholars have identified significant doctrinal vulnerabilities in Jackson's taxonomy. Mark D. Rosen argued in the \textit{UCLA Law Review} that Jackson's Zone 3 embeds an unexamined assumption of ``categorical congressional supremacy'' whenever congressional and presidential powers overlap, without adequately justifying why Congress always trumps the Commander in Chief in shared domains (\href{https://www.uclalawreview.org/wp-content/uploads/2019/09/44_54UCLALRev1703August2007.pdf}{UCLA Law Review}). More recently, scholars analyzing \textit{Trump v. United States}, 603 U.S. 593 (2024), have warned that expanding Zone 3 ``exclusive'' presidential powers to confer broad criminal immunity risks inverting Jackson's core anti-aggrandizement premise (\href{https://pennlawreview.com/2026/03/06/power-and-immunity-in-youngstown-and-trump-v-united-states/}{Penn Law Review}).

\textbf{Underlying Competing Values:} \textit{Youngstown} embodies the central structural tension of American constitutionalism: balancing the need for an energetic, decisive Executive capable of protecting national security during unforeseen crises against the Founders' conviction that diffusing lawmaking authority in a deliberative legislature is essential to prevent arbitrary power and preserve the rule of law.

\section{Key Quotations}

\begin{quote}
``The President's power, if any, to issue the order must stem either from an act of Congress or from the Constitution itself.''
\end{quote}
\noindent --- Justice Black, Opinion of the Court, 343 U.S. at 585 (\href{https://supreme.justia.com/cases/federal/us/343/579/}{Justia}; \href{https://www.courtlistener.com/opinion/105018/youngstown-sheet-tube-co-v-sawyer/}{CourtListener}).

\begin{quote}
``In the framework of our Constitution, the President's power to see that the laws are faithfully executed refutes the idea that he is to be a lawmaker. The Constitution limits his functions in the lawmaking process to the recommending of laws he thinks wise and the vetoing of laws he thinks bad.''
\end{quote}
\noindent --- Justice Black, Opinion of the Court, 343 U.S. at 587--588 (\href{https://supreme.justia.com/cases/federal/us/343/579/}{Justia}).

\begin{quote}
``In short, a systematic, unbroken, executive practice, long pursued to the knowledge of the Congress and never before questioned, engaged in by Presidents who have also sworn to uphold the Constitution, making as it were such exercise of power part of the structure of our government, may be treated as a gloss on `executive Power' vested in the President by \S{} 1 of Art. II.''
\end{quote}
\noindent --- Justice Frankfurter, concurring, 343 U.S. at 610--611 (\href{https://supreme.justia.com/cases/federal/us/343/579/}{Justia}).

\begin{quote}
``While the Constitution diffuses power the better to secure liberty, it also contemplates that practice will integrate the dispersed powers into a workable government. It enjoins upon its branches separateness but interdependence, autonomy but reciprocity. Presidential powers are not fixed but fluctuate depending upon their disjunction or conjunction with those of Congress.''
\end{quote}
\noindent --- Justice Jackson, concurring, 343 U.S. at 635 (\href{https://supreme.justia.com/cases/federal/us/343/579/}{Justia}; \href{https://tile.loc.gov/storage-services/service/ll/usrep/usrep343/usrep343579/usrep343579.pdf}{Library of Congress}).

\begin{quote}
``Under this messenger-boy concept of the Office, the President cannot even act to preserve legislative programs from destruction so that Congress will have something left to act upon.''
\end{quote}
\noindent --- Chief Justice Vinson, dissenting, 343 U.S. at 708--709 (\href{https://supreme.justia.com/cases/federal/us/343/579/}{Justia}; \href{https://static.c-span.org/landmarkCases/pdf/Youngstown_Vinson_Dissent.pdf}{C-SPAN Vinson Dissent}).

\end{document}
